\section{Polynomial occupation cuts and the full fixed-band spectrum}
\label{sec:v40-occupation-powers}
We estimate the true section occupation at every real occupation
frequency. The extra discontinuities are counted on an already regular
inverse collision curve, not by multiplying the number of cells in a
one-step partition. This distinction changes the complexity bound from
an exponential bound to a polynomial one.

Throughout this section the table is circular, $R\in I$, and the
section is the original $Y_R^*$. Put
\[
 \mathscr T_{R,\theta}=e^{i(b\bar\tau_R+vc)}\mathcal Q_{R,(u,b),v},
 \qquad \theta=(u,b,v)\in\T^2\times\R\times\T.
\]
This is the uncentered operator. On physical densities its action is
\begin{equation}\label{eq:v40-physical-operator}
 \mathscr T_{R,\theta}h
   =\bigl(e^{i(u\cdot\kappa_R+b\tau_R+v\eta_R)}h\bigr)\circ T_R^{-1}.
\end{equation}
It is defined on the local strong and weak spaces by the preceding
flight and section multipliers already proved. The expression does
not require the forward roof alone to be a strong-space multiplier.
The torus convention is $\T=\R/(2\pi\Z)$; its metric is shortest
angular distance, combined with Euclidean distance in the real roof
coordinate. The roof band $B$ is always fixed before any limit in the
collision count.

\subsection{A refinement of the existing inverse-curve partition}
Retain $\mathcal G_m(W),\mathcal I_m(W),j_U,\varsigma$ and the four
geometric sums of Section~\ref{sec:action-norm-details}.

\begin{lemma}[Linear occupation complexity on one inverse curve]
\label{lem:v40-linear-cuts}
There is $K<\infty$, independent of $m,R,W$, such that every
$U\in\mathcal G_m(W)$ can be cut into at most $r_m=1+Km$ intervals
on which $(\eta_R\circ T_R^j)_{0\le j<m}$ is constant. On each such
interval the accumulated uncentered phase has $C^p$ norm at most
$C_B$, for $|b|\le B$, uniformly in $u,v,m,R$.

If $\mathcal G_m^\eta(W)$ denotes this refinement, and $j_V$ is the
supremum of the same stable Jacobian restricted to $V$, then
\begin{align}
 \sum_{V\in\mathcal G_m^\eta(W)}j_V&\le Cr_m,
                                      \label{eq:v40-cut-sums-a}\\
 \sum_{V\in\mathcal G_m^\eta(W)}(|V|/|W|)^\varsigma j_V&\le Cr_m,
                                      \label{eq:v40-cut-sums-b}\\
 \sum_{\substack{U\in\mathcal I_m(W)\\V\subset U}}
            (|V|/|W|)^\varsigma j_V
       &\le Cr_m\theta_*^{(1-\varsigma)m},
                                      \label{eq:v40-cut-sums-c}\\
 \sum_{V\in\mathcal G_m^\eta(W)}j_V^{1-\varsigma}&\le Cr_m C_3^m.
                                      \label{eq:v40-cut-sums-d}
\end{align}
The designation ``without a long ancestor'' in the third sum is that
of the original collision tree; no new growth lemma for an iterated
section partition is being assumed.
\end{lemma}
\begin{proof}
Every intermediate image $T_R^jU$ is a regular stable graph. Its slope
is negative and bounded away from zero and infinity. A horizontal or
vertical side of one of the five section rectangles is therefore
crossed at most once, after the fixed angular chart cuts. This is the
same transversality used in Lemma~\ref{lem:v37-section-multiplier}.
There are a fixed number of sides. Each of the $m$ intermediate images
contributes at most that number of cut points on $U$, since the map
from $U$ to its image is one-to-one. Their union has at most $Km$
points, not $K^m$ points.

On a component the occupation phase is constant, of modulus one.
The displacement label is also constant on the old regular branch.
The roof action estimate \eqref{eq:v35-action-telescope} supplies the
remaining $C^p$ bound. Normalized stable-Jacobian distortion bounds
are unchanged by restriction to a subinterval.

For nonnegative lengths $l_1,\ldots,l_r$ with sum $l$,
$\sum l_i^\varsigma\le r^{1-\varsigma}l^\varsigma$. Apply this to the
pieces of each $U$, and use $j_V\le j_U$. The first and fourth sums
cost at most $r_m$ times their old sums. The second and third cost at
most $r_m^{1-\varsigma}$ times the corresponding old length sums.
The weaker displayed bounds follow since $r_m\ge1$.
\end{proof}

We also need compatibility with the unstable comparison; counting
cuts on a single curve would not by itself establish this.

\begin{lemma}[Occupation matching and its unmatched cost]
\label{lem:v40-occupation-matching}
The matched pairs of Section~\ref{sec:action-norm-details} admit a
refinement with at most $Cr_m$ pieces per old pair. On each retained
pair the occupation itineraries agree. Their graph distance retains
the old $C\Lambda^{-m}\epsilon$ bound, and the two accumulated weights
have $C^q$ difference at most $C_B\epsilon^{1-q}$.
The additional unmatched pieces $V$ satisfy
\begin{equation}\label{eq:v40-new-unmatched}
 |T_R^mV|\le C\epsilon,\qquad
 \sum_V j_V^{1-\varsigma}\le Cr_m C_3^m.
\end{equation}
All constants are uniform in the occupation frequency.
\end{lemma}
\begin{proof}
Start with the regular connector foliation of an old matched pair.
At time $j$, its connectors have length at most
$C\Lambda^{-(m-j)}\epsilon$, by backward contraction of unstable
connectors. Remove connectors meeting a section side. Transversality
of that side with the two stable graphs bounds the length removed on
either graph at time $j$ by $C\Lambda^{-(m-j)}\epsilon$ per crossing.
Forward stable contraction from $j$ to $m$ makes the corresponding
terminal length at most $C\Lambda^{-2(m-j)}\epsilon\le C\epsilon$.
There are at most a fixed number of crossings at each time, by the
preceding lemma.

Use closed removed intervals before trimming to a common graph domain.
If changing from the connector pairing to the common graph-coordinate
pairing would cross an occupation cut, remove that interval as well.
The old trimming moves a point by at most $C\epsilon$ along an inverse
stable graph. At time $j$ this displacement is at most
$C\Lambda^{-j}\epsilon$; its terminal image is at most
$C\Lambda^{-m}\epsilon$. It has the same admissible unmatched cost.
This construction has at most $Cr_m$ endpoints per old pair and
ensures agreement of the entire half-open occupation itinerary on
the retained common domains. Restriction to equal graph domains
preserves the old bound on the distance of the inverse graphs.

On retained pairs the occupation factors cancel exactly in the weight
difference. The old roof-action estimate and interpolation give
$C_B\epsilon^{1-q}$, without an extra factor growing with $m$.
For every removed piece bounded distortion gives
$j_V|V|\le C\epsilon$. The last sum in
\eqref{eq:v40-new-unmatched} follows by charging each restriction to
its old parent and using the $Cr_m$ multiplicity and
\eqref{eq:v35-geometric-sums-b}. Thus their stable-norm contributions
sum to at most
$C_B r_m C_3^m\epsilon^\varsigma\|h\|_s$.
The old collision singularities and homogeneity cuts remain in the
old unmatched family throughout. No connector through such a cut is
used in the argument.
\end{proof}

\subsection{The three estimates and their spectral consequence}
\begin{theorem}[Full occupation-torus inequalities]
\label{thm:v40-polynomial-LY}
For each fixed $B<\infty$, all real $\theta$ with $|b|\le B$, and
$m\ge0$, the local common collision norms satisfy
\begin{align}
 |\mathscr T_{R,\theta}^m h|_w
     &\le C_B(1+m)|h|_w,                        \label{eq:v40-full-LY-w}\\
 \|\mathscr T_{R,\theta}^m h\|_s
     &\le C_B(1+m)(\theta_*^{(1-\varsigma)m}+\Lambda^{-qm})\|h\|_s
                +C_B(1+m)|h|_w,               \label{eq:v40-full-LY-s}\\
 \|\mathscr T_{R,\theta}^m h\|_u
     &\le C_B(1+m)\Lambda^{-\gamma m}\|h\|_u
                +C_B(1+m)C_4^m\|h\|_s.        \label{eq:v40-full-LY-u}
\end{align}
In a $B$-dependent equivalent strong norm there are $\sigma_B<1$ and
$C_B'<\infty$ such that
\begin{equation}\label{eq:v40-full-DF}
 \|\mathscr T_{R,\theta}^m h\|_B
       \le C_B'\sigma_B^m\|h\|_B+C_B'(1+m)|h|_w.
\end{equation}
The essential spectral radius is uniformly less than one on the full
fixed band and occupation torus. The operators vary locally
strong-to-weak continuously in $(R,\theta)$.
\end{theorem}
\begin{proof}
For the weak norm split the weighted curve pairing
\eqref{eq:weighted-curve-pairing} over the intervals in
Lemma~\ref{lem:v40-linear-cuts}. Each normalized test has $C^p$ norm
at most $C_B$. Equation~\eqref{eq:v40-cut-sums-a} proves
\eqref{eq:v40-full-LY-w}.

For the strong stable norm subtract the average $a_U$ of the pulled-back
terminal test on each \emph{old} inverse piece $U$. Keep that same
average on every occupation subinterval $V\subset U$. Stable
contraction gives
$\|\psi\circ T^m-a_U\|_{C^q(U)}\le
 C\Lambda^{-qm}|W|^{-\varsigma}$.
Its refined sum is controlled by \eqref{eq:v40-cut-sums-b}.
The average on old never-long pieces is controlled by
\eqref{eq:v40-cut-sums-c}. For all other pieces retain their old last
long ancestor. The weak payment of each descendant now has at most
$r_m$ summands. In the notation of Section~\ref{sec:action-norm-details}
its total is bounded by
\[
 C_B r_m\delta_0^{-\varsigma}|h|_w
 \sum_{k=0}^m\theta_*^{m-k}
 \sum_{V\in\mathcal G_k(W)}(|V|/|W|)^\varsigma j_V
       \le C_B r_m|h|_w.
\]
This is the last-long sum from \cite[equations (4.7)--(4.11)]{DZ},
with its length factors unchanged and an explicitly counted
multiplicity. Newly short occupation pieces are not incorrectly put
in the old never-long family. This proves \eqref{eq:v40-full-LY-s}.

For the unstable norm use Lemma~\ref{lem:v40-occupation-matching}.
The old matched leading comparison acquires at most $Cr_m$ copies.
The unmatched contribution is bounded there. The other errors have
powers $\epsilon^{p-q},\epsilon^{1/3-q},\epsilon^{1-q}$ and the old
Jacobian sum, now with multiplicity at most $Cr_m$. Each of these
powers, as well as $\varsigma$, exceeds $\gamma$ by
\eqref{eq:v35-exponent-choice}. Division by $\epsilon^\gamma$
proves \eqref{eq:v40-full-LY-u}, after enlarging $C_4$. This is the
matched/unmatched calculation of \cite[Section 4.3]{DZ}, not an
inference from the one-step multiplier norm.

Choose $N=N(B)$ so that both leading strong coefficients at time $N$
are below $1/4$. This is possible despite their factor $1+N$. Next
choose $a_B>0$ so that the stable cross term, multiplied by $a_B$,
is below $1/4$. With $\|h\|_B=\|h\|_s+a_B\|h\|_u$ this gives
$\|\mathscr T^N h\|_B\le\frac12\|h\|_B+D_B|h|_w$.
Iteration uses the already established bound
$|\mathscr T^{jN}h|_w\le C_B(1+jN)|h|_w$.
The geometric convolution with $2^{-j}$ is at most $C_B(1+m)$;
treating the remaining powers below $N$ proves
\eqref{eq:v40-full-DF}. Compact embedding gives the essential-radius
assertion. This conclusion needs no claim that the weak powers are
uniformly bounded.

For parameter continuity, the roof operator is already locally
strong-to-weak continuous. The symmetric difference of the two moving
sections lies in finitely many strips whose width tends to zero.
Its intersection with a stable curve consists of a bounded number of
intervals of length $O(\delta)$, where $\delta$ bounds the motion of
the rectangle sides. The strong stable norm bounds their weak pairing
by $C\delta^\varsigma\|h\|_s$. Thus
$\|M_{\eta_R}-M_{\eta_{R_0}}\|_{\mathcal B\to\mathcal B_w}
\le C\delta^\varsigma$. The exact affine formula
\eqref{eq:v37-occupation-exponential}, weak boundedness of one roof
step, and the existing roof comparison give the assertion on the full
occupation torus. Frequency continuity is strong analyticity in each
local coordinate chart. Completion extends the estimates from smooth
physical densities to the stated spaces.
\end{proof}

\begin{theorem}[Physical peripheral spectrum on the full fixed band]
\label{thm:v40-full-peripheral}
For every fixed $B$, $\sup_{R,\theta,m}\|\mathscr T_{R,\theta}^m\|_B$
is finite for $|b|\le B$. At a frequency $\theta$ there is a
unit-modulus eigenvalue if and only if $\theta\in\mathcal G_R$.
It is then the single simple eigenvalue $e^{-is_R(\theta)}$, with a
bounded physical eigenvector $q_\theta\nu$ satisfying
\eqref{eq:v39-phase-equation}.

Every compact subset of
$\{(R,\theta):|b|\le B,\ \theta\notin\mathcal G_R\}$ admits
constants $C<\infty$ and $\rho<1$ for which
$\|\mathscr T_{R,\theta}^m\|_B\le C\rho^m$.
This statement is about a compact subset of the \emph{joint}
parameter-frequency complement, not a pointwise choice of gaps.
\end{theorem}
\begin{proof}
The polynomial power bound in \eqref{eq:v40-full-DF} excludes spectral
radius greater than one. Quasi-compactness leaves a finite-dimensional
peripheral space at each fixed parameter and frequency. We first
exclude its Jordan blocks, without presupposing power boundedness.
For a smooth bounded initial density and a smooth global test the
physical pairing of every power is bounded independently of $m$.
Its spectral decomposition is a finite sum of exponential polynomials
plus an exponentially decaying term. A bounded exponential polynomial
has no positive-degree coefficient: divide by its highest power of
$m$, multiply by a conjugate peripheral phase, and take Cesaro means.
Distinct phases have vanishing cross means. Faithfulness on the strong
space, Lemma~\ref{lem:v35-faithful-strong}, therefore makes every
positive-degree vector coefficient zero. Density of smooth sources
in the strong space eliminates every nonzero nilpotent peripheral
part. The powers at this fixed point are now bounded.

The physical Cesaro argument of
Lemma~\ref{lem:v35-peripheral-density} applies verbatim to
\eqref{eq:v40-physical-operator}. Each peripheral vector has a bounded
physical representative. The modulus is constant by ergodicity.
Quotients of two such circle-valued representatives solve an ordinary
collision eigenfunction equation. Mixing makes their eigenvalues
equal, and ergodicity makes their quotient constant. Thus the
peripheral eigenvalue is unique and simple. Its physical equation is
exactly \eqref{eq:v39-phase-equation} with eigenvalue $e^{-is}$.

Conversely, if a measurable phase exists and all strong eigenvalues
are strictly inside the unit circle, smooth endpoint pairings decay
exponentially. Approximate the measurable phase in $L^1$, keep that
approximation fixed while $m\to\infty$, and then remove the
approximation, just as in Lemma~\ref{lem:v39-local-phase-exclusion}.
Its unit-modulus pairing gives a contradiction. This proves the
converse without assuming that an arbitrary measurable multiplier
acts on the strong space.

Finally apply the strong-to-weak perturbation theorem with
\eqref{eq:v40-full-DF}; the factor $1+m$ is bounded by $C_\epsilon
(1+\epsilon)^m$. Near a point without peripheral spectrum all powers
decay on a common contour. Near a point with peripheral spectrum the
single rank-one projection persists, the complementary powers decay
on a common contour, and the remaining eigenvalue has modulus at
most one by the global polynomial bound. Its projection is locally
uniformly bounded. A finite cover of the compact parameter-frequency
band gives uniform power boundedness. The same argument on a compact
joint nonresonant subset gives a uniform strictly smaller radius.
It also shows that the joint resonance locus is closed. Constants
may depend on $B$; no growing-band estimate is asserted.
\end{proof}
